\documentclass{article}
\usepackage{amsmath}
\usepackage{geometry}
\usepackage[colorlinks=true, linkcolor=blue, urlcolor=blue]{hyperref}
\usepackage{enumitem}
\usepackage{csquotes}
\geometry{a4paper, margin=1in}
\usepackage[style=apa,backend=biber]{biblatex}
\addbibresource{sources.bib}

\title{An Extensible Benchmarking Framework for Federated
Synthetic Tabular Data Generators}
\author{Thomas Stavik Lønning \and Alexander Benjamin Rød Opedal \and Ken Solbakken Remen \and Trond Thorsteinsson \and Bård Linga}
\date{\today}

\begin{document}
\raggedbottom
\maketitle
\tableofcontents
\newpage

\section{Introduction}
Hello FedBench! Have you been BenchFed?

\section{Background}

\subsection{Federated Learning}
In machine learning, the main approach that has been used in the past
is to collect data on a central server, and then use machine learning
algorithms to train some model on the data. In Federated Learning (FL),
rather than moving the data to the training, the training
is moved to the data \parencite{flower_tutorial}.
In a typical client-server architecture, locally trained model updates
are then sent back to the server for aggregation, and this is repeated
until the global model converges. FL thus enhances privacy by keeping
potentially sensitive data local to the owning device.

\subsection{The Flower framework}
In the words of \cite{fed_frameworks}, \enquote{a FL framework provides the
infrastructure support to build collaborative machine learning
models across multiple devices, servers, or entities without
transferring raw data to a central location. Instead, the models
are trained locally on individual devices, servers or entities
and only the model updates are exchanged and aggregated to
create a global model or centralized model.}

Flower \parencite{flower_article} is a FL framework written in python, and designed to be
scalable, client-agnostic, communication-agnostic, privacy-agnostic
and flexible. In Flower, global logic, such
as client selection, configuration, parameter update aggregation,
and federated or centralized model evaluation is expressed through
Flower's \textit{Strategy} abstraction. A concrete \textit{Strategy}
implementation represents a single FL algorithm, and Flower provides
tested reference implementations of many popular algorithms.
\parencite{flower_article}

Local logic can be expressed using the high-level \textit{Client}
abstraction, but it is also possible to integrate with Flower at the
protocol level, meaning you can in theory implement clients in
virtually any programming language.
\parencite{flower_article}

Built into Flower is a Virtual Client Engine (VCE), a tool that
enables virtualization of Flower Clients, which in turn enables
large scale single-machine or multiple-machine experiments.
\parencite{flower_article}

\subsection{Evaluation Metrics}
The use of synthetic data in high-risk AI applications related to
human health warrants rigorous data quality assessment. \cite{metrics}
conducted a comprehensive literature review on the quality evaluation
metrics within the scope of synthetic tabular healthcare data, and
highlighted the need for objective evaluation standards.
They define five relevant quality dimensions:
\begin{itemize}[label=--]
  \item \textbf{Similarity}: The similarity between the synthetic and
    real dataset, and the ability to mimic statistical characteristics and
    patterns of the original dataset.
  \item \textbf{Usability}: The usability of inferences drawn from
    the synthetic data.
  \item \textbf{Privacy}: The risk of re-identification for individuals
    from the original dataset.
  \item \textbf{Fairness}: Underrepresentation of certain groups that
    may lead to unfair use. Transparency of the generation process and output.
  \item \textbf{Carbon footprint and computational complexity}: Resource
    use for preprocessing, model training, and inference.
\end{itemize}

A critical goal of this thesis is to provide concrete, reusable implementations
covering most, if not all of these dimensions, and thus provide a basis for
standardized evaluation of federated synthetic data generators.

\section{Project description and objectives}

\subsection{Background and motivation}

\subsection{Problem statement}

\subsection{Objectives and research questions}

\subsection{Scope, assumptions, and limitations}

\section{Related work}

\subsection{Centralized synthetic tabular data generation}

\subsection{Federated learning for tabular data}

\subsection{Federated synthetic tabular data generation}

\subsection{How synthetic data is evaluated}

\section{System design: FedBench}

\subsection{Design goals}

\subsection{Architecture overview}

\subsection{Core modules}
\subsubsection{Orchestration layer (Flower)}
\subsubsection{Data layer (loading, preprocessing, partitioning)}
\subsubsection{Generator layer (interfaces and implementations)}
\subsubsection{Evaluation layer (metrics, attacks, and reporting)}
\subsubsection{Experiment configuration, logging, and reproducibility}

\subsection{Extensibility and API design}
\subsubsection{Abstract base classes and required methods}
\subsubsection{Registration and discovery of generators}

\section{Data}

\subsection{Datasets}
\subsubsection{UCI Heart Disease}
\subsubsection{Breast Cancer Wisconsin}
\subsubsection{NCCTG Lung Cancer}

\subsection{Preprocessing and feature handling}

\subsection{Federated partitioning schemes}
\subsubsection{Horizontal partitioning (Fed-TGAN)}
\subsubsection{Vertical partitioning (VT-GAN)}

\section{Methods}

\subsection{Federated synthetic data generators}
\subsubsection{Fed-TGAN}
\subsubsection{VT-GAN (GTV)}
\subsubsection{FedTabDiff}

\subsection{Experimental protocol}
\subsubsection{Federated learning setup (clients, rounds, aggregation)}
\subsubsection{Hyperparameters and tuning strategy}
\subsubsection{Compute environment}

\subsection{Evaluation methodology}
\subsubsection{Fidelity and utility (TSTR, similarity, correlations)}
\subsubsection{Privacy (membership/attribute inference, DP accounting)}
\subsubsection{Fairness (demographic parity, equalized odds)}
\subsubsection{Scalability (runtime, communication overhead, compute footprint)}

\section{Implementation}

\subsection{Framework implementation details}

\subsection{Integration of reference generators}

\subsection{How to extend the framework (short guide)}

\section{Results}

\subsection{Fidelity and utility results}

\subsection{Privacy results}

\subsection{Fairness results}

\subsection{Scalability results}

\section{Discussion}

\subsection{Summary of findings}

\subsection{Trade-offs and recommendations}

\subsection{Threats to validity}

\section{Conclusion and future work}

\subsection{Conclusion}

\subsection{Future work}

\appendix
\section{Appendix}
\subsection{Reproducibility checklist}

\printbibliography

\end{document}
